\section{Pre-evaluation research state}
\label{sec:pre-evaluation}
Deze slotstatus begrenst precies wat theoretisch wordt voorgesteld en wat nog onderzocht moet worden. De centrale these T betreft praktisch relevante nettowaarde van contextbinding voor de vooraf afgebakende identiteit-, tijd- en gezagstaken. T beweert niet dat iedere organisatie een centrale ontologie of een RDF-database nodig heeft. De gerichte bronreview maakt de novelty-vraag scherper, maar beantwoordt haar niet definitief.

\subsection{Theoretische propositions en falsificatie}
\begin{longtable}{@{}P{.08\linewidth}P{.42\linewidth}P{.42\linewidth}@{}}
\caption{Niet-getoetste propositions en hun weerleggingscondities.}\\
\toprule ID & Kandidaatpropositie en grondslag & Ondersteuning versus tegenbewijs \\ \midrule\endfirsthead
\toprule ID & Kandidaatpropositie en grondslag & Ondersteuning versus tegenbewijs \\ \midrule\endhead
TP1 & Expliciete contextbinding kan bij voldoende bronkwaliteit interpretatie verbeteren boven dezelfde informatie in A1. Rationale: context- en metadatarepresentatie. \citep{S08,S33,L13} & E1: relevante winst met kosten-/foutgrenzen ondersteunt H1. Praktische gelijkwaardigheid of nadeel met voldoende precisie weerspreekt meerwaarde; extra informatie alleen ondersteunt het mechanisme niet. \\
TP2 & Betrouwbare epistemische annotatie kan onterechte gezagsopwaardering verminderen; misleidende annotatie kan dit voordeel omkeren. Eigen synthese uit provenance en gezagscontext. \citep{S33,S25} & E2: minder fouten bij gelijke inhoud ondersteunt de kandidaat. Geen relevante afname, meer fouten of uitsluitend labelblind volgen vormt tegenbewijs voor de gespecificeerde gebruiksconditie. \\
TP3 & Ontologische verdieping kan alleen netto waarde toevoegen bij relevante identiteit-/rolproblemen en aanvaardbare extra inspanning. \citep{L08,L09,L10} & E3: betere foutdetectie binnen tijdsgrens ondersteunt gerichte inzet. Gelijke detectie, meer fout-positieven of disproportionele kosten weerlegt de gekozen meerwaardeclaim. \\
TP4 & Versiegebonden afhankelijkheden kunnen reconstructie/impact verbeteren, mits mappings en contextresolutie voldoende betrouwbaar zijn. \citep{L14,S33,S47} & E4: vooraf relevante winst ondersteunt de kandidaat. Gelijkwaardige archiefoplossing, slechter herstel of meer fouten weerlegt het gekozen voordeel; onzekerheid blijft onbeslist. \\
\bottomrule\end{longtable}

De propositions zijn conditioneel; hun toetsbare vorm vereist vooraf afgebakende populaties en marges. De woorden ``kan'' of ``voldoende'' mogen niet na een negatieve uitkomst worden gebruikt om de taakpopulatie opportunistisch te verplaatsen. Wijziging van voorwaarden na een test creëert een nieuwe hypothese. TP2--TP4 kunnen afzonderlijk worden verworpen zonder dat onterecht een universele uitspraak over alle semantiek wordt gedaan.

\subsection{Afgeleide principes en te onderzoeken prototype-eigenschappen}
DP1--DP9 blijven kandidaatprincipes. Zij betreffen respectievelijk getypeerde rollen en mappings, uitspraakgebonden gezag, identiteit/versie, historie, scheiding van inferentie en kwaliteit, productcontract, beheerste verandering, selectieve ontologische verdieping en epistemische metadata. Hun bron-naar-ontwerpafleiding staat in het derivation-register; geen principe heeft de status empirisch gevalideerd.

Het prototype moet nog feitelijk worden geïdentificeerd met repository/release, configuratie en gegevens-/modelversies. Daarna moeten minimaal worden onderzocht: onderscheiden representatierollen (R01/R04), bron- en gezagsclaims (R02/R08/R23), identifiergedrag (R03), decoderings- en modelcontracten (R05/R11/R14), historie en provenance (R06/R07/R20), toetsbare kwaliteit (R09/R16), release/lifecycle en events (R10/R12/R15/R19), product-/toepasselijkheidsmetadata (R13/R18/R21), ontologische opties (R22) en consumentgedrag (R17). De huidige status van elk van deze eigenschappen is \emph{niet beoordeeld}; afwezigheid van evidence in dit dossier betekent niet dat de functie ontbreekt.

\subsection{Benodigde experimenten en open beslissingen}
E1--E4 in tabel~\ref{tab:experiments} specificeren de onafhankelijke en afhankelijke variabelen. Representatie, annotatievorm/-kwaliteit, modelreviewmethode en afhankelijkheids-/resolutieconditie zijn mogelijke interventies. Correctheid, herleidbaarheid, onterechte opwaardering, foutdetectie, reconstructie, tijd en kosten zijn onderscheiden uitkomsten. Menselijke expertise, taaktype en AI-modelversie worden als gecontroleerde of gestratificeerde factoren behandeld; er zijn nog geen data verzameld.

Vóór uitvoering ontbreken nog volledige relevante normteksten, een onafhankelijke inhoudelijke review, een gecontroleerde meerregistratiecasus, een geadjudiceerde gouden standaard, een steekproefonderbouwing en vastgelegde relevantie-/acceptatiemarges. Het onderzoeksprotocol verbiedt het laten vallen van mislukte taken omdat de context ontbreekt. Een eenvoudige A1-oplossing die gelijkwaardig of beter presteert tegen lagere kosten moet tot verwerping of vereenvoudiging van het rijkere ontwerp kunnen leiden. Dat is een beoogde wetenschappelijke uitkomst, geen ongewenste uitzondering op het onderzoeksplan.
